\documentclass[a4paper,oneside,11pt]{amsart}

%\usepackage[spanish, activeacute]{babel}
\usepackage[utf8]{inputenc}
%\usepackage[latin1]{inputenc}
\usepackage{latexsym}
\usepackage{multicol}
\usepackage{enumerate}
\usepackage{enumitem}
\usepackage{booktabs}
\usepackage[heightrounded]{geometry}
\geometry{top=2cm,bottom=2cm,left=2cm,right=2cm}
\usepackage{graphicx}

\DeclareMathOperator{\cond}{cond}
\DeclareMathOperator{\Id}{Id}
\DeclareMathOperator{\re}{Re}
%\DeclareMathOperator{\im}{Im}
\newcommand{\I}{\mathbb{I}}
\newcommand{\R}{\mathbb{R}}
\newcommand{\Q}{\mathbb{Q}}
\newcommand{\C}{\mathbb{C}}
\newcommand{\Z}{\mathbb{Z}}
\newcommand{\N}{\mathbb{N}}
\DeclareMathOperator{\im}{Im}

\newcommand{\nombremateria}
{Investigaci\'on Operativa}
\newcommand{\nombrecuatrimestre}
    {Segundo cuatrimestre de 2022}
\newcommand{\numeroparcial}
    { Segundo Parcial }
\newcommand{\fecha}
    {11 de Octubre de 2022}
\newcommand{\numerotema}
    {\bf 1}

%\oddsidemargin -1cm \evensidemargin -1cm \textwidth 17.5cm
%topmargin 1.2cm \textheight 25cm
%\setlength{\textheight}{25cm}

\def \ds{\displaystyle}
\theoremstyle{definition}
\newtheorem{quest}{Ejercicio}
\thispagestyle{empty}

\begin{document}

\hfill \hfill
\begin{tabular}[c]{|c|c|c|c|c|}
\hline 
1 & 2 & 3 & 4 & Calificaci\'on \\ \hline \hline
\hspace{1 cm} & \hspace{1 cm} & \hspace{1 cm} & \hspace{1 cm} & \hspace{1.3 cm} \\
\hspace{1 cm} & \hspace{1  cm} & \hspace{1  cm} & \hspace{1  cm} & \hspace{1.3 cm} \\ \hline

\end{tabular}

\vspace{-1cm}

\noindent Nombre y apellido:

\vspace{0.2cm}

\noindent Número de libreta:

\noindent \hrulefill 

\smallskip
\renewcommand{\baselinestretch}{1.1}

\begin{center}
	\large \textbf{\nombremateria} 
	
	\numeroparcial -- \fecha
\end{center}

\noindent\textit{Para aprobar el parcial deberá obtenerse un puntaje mayor o igual que 4. \textbf{Entregue todos los ejercicios en hojas separadas.}}
% \vspace{.1cm}
\renewcommand{\theenumi}{\textbf{\arabic{enumi}}}

%%% EMPIEZA EL EJERCICIO 1 %%%
\fontsize{10.5pt}{0.4cm}\selectfont
\begin{quest} \verb|[4 pts.]| Considerar el siguiente modelo de Programación Lineal:

\[
\begin{array}{rrcl}
\max & \multicolumn{3}{l}{z = -x_1+2x_2-x_3+x_4} \\
s.a: & -x_1-x_3 & \leq & -2 \\
 & 3x_2-x_3+x_4 & \leq & 6 \\
 & x_1+2x_2-x_4 & = & 10 \\
 & x_1,x_2,x_4 & \geq & 0 \\
 & x_3 & \leq & 0 \\
 
\end{array}
\]
\begin{enumerate}[label=\alph*)]
    \item \verb|[2.5 pts.]| Hallar una solución óptima del problema mediante el Método de Dos Fases.
    \item \verb|[1.5 pt.]| Plantear el problema dual y hallar su óptimo.
\end{enumerate}

\end{quest}

%%% TERMINA EL EJERCICIO 1 %%%

%%% EMPIEZA EL EJERCICIO 2 %%%

\begin{quest} \verb|[1 pt.]| El siguiente diccionario corresponde a un problema de Programación Lineal con objetivo de \textbf{minimizar}:
    \[
    \begin{array}{rrrrrrrrr}
        x_3 & = & 1   & - & 2x_1 & + & 2x_5 & - & x_6   \\
        x_2 & = & 1+a & - & bx_1 & - &  x_5 & - & 2x_6  \\
        x_4 & = & 1   & + & x_1  & + & x_5 & - & x_6   \\ \hline
        z   & = & -1  & + & cx_1 &   &      & + & 2ax_6  \\
    \end{array}
    \]
    con $a,b,c\in\mathbb{R}$. Utilizando el criterio usual de elección de qué variable entra a la base, decidir si cada una de las siguientes afirmaciones son verdaderas o falsas y justificar la respuesta.
    \begin{enumerate}[label=\alph*)]
    \item \verb|[0.25 pts.]| Si $c> 0$ y $a> 0$, la base actual es la única base óptima.
    \item \verb|[0.25 pts.]| Si $a=-1$ y $c< -2$, la próxima iteración es degenerada.
    \item \verb|[0.25 pts.]| Si $c < -2$, $x_1$ entra a la base en la siguiente iteración.
    \item \verb|[0.25 pts.]| Si $a=1$, $c=0$ y $b=-2$, entonces $(\frac{1}{4}, \frac{5}{2}, \frac{1}{2},\frac{5}{4}, 0, 0)$ es una solución óptima.  
    \end{enumerate}

\end{quest}
%%% TERMINA EL EJERCICIO 2 %%%

%%% EMPIEZA EL EJERCICIO 3 %%%


\begin{quest} \verb|[2.5 pts.]| Para $\alpha=\beta=\gamma=0$, se sabe que una solución óptima del siguiente problema es $x^\ast = (4,2,0)$:

\[
\begin{array}{rrcl}
\max & \multicolumn{3}{l}{3x_1+(2+\beta)x_2+(1+\gamma)x_3} \\
s.a: & x_1+2x_2+x_3 & \leq & 8+\alpha \\
 & -x_1+x_2-x_3 & \leq & 6 \\
 & 2x_1+x_3 & \leq & 8 \\
 & x_1, x_2, x_3 & \geq & 0
\end{array}
\]

    \begin{enumerate}[label=\alph*)]
    \item \verb|[1 pt.]| Fijando $\beta=\gamma=0$, hallar todos los $\alpha\in\mathbb{R}$ para los cuales la base óptima siga siendo factible. Calcular el valor de $x^\ast$ para $\alpha=2$. 
    \item \verb|[1.5 pts.]| Fijando $\alpha=0$, hallar un conjunto $S\subseteq \mathbb{R}^2$ tal que si $(\beta,\gamma)\in S$, entonces $x^\ast$ sigue siendo solución óptima.
    \end{enumerate}

\end{quest}

%%% TERMINA EL EJERCICIO 3 %%%

%%% EMPIEZA EL EJERCICIO 4 %%%


\begin{quest} \verb|[2.5 pts.]| Utilizar el algoritmo de Branch \& Bound para resolver el siguiente problema de Programación Lineal Entera, detallando cada división en subproblemas mediante el esquema de árbol visto en clase. Comenzar ramificando por $x_2$.
\[
\begin{array}{rrcl}
\max & \multicolumn{3}{l}{3x_1+2x_2} \\
s.a: & -2x_1+x_2 & \leq & 4 \\
 & 2x_1+x_2 & \geq & 4 \\
 & 2x_1+4x_2 & \leq & 17 \\
 & -3x_1+3x_2 & \geq & -11 \\
 & x_1, x_2 & \geq & 0 \\
 & x_1, x_2 & \in & \mathbb{Z}
\end{array}
\]
\end{quest}
%%% TERMINA EL EJERCICIO 4 %%%

%%% TERMINA EL PARCIAL %%%



\end{document}

